\subsection{内射余生成元}

命题11.——设B是一个环，F是一个B-模，P是一个(B, A)-双模。若F是内射B-模且P是平坦A-模， \(\mathrm{Hom}_{\mathbf{B}}(\mathbf{P},\mathbf{F})\) 是内射A-模。

设 \(u: \mathbf{M}' \to \mathbf{M}\) 是A-模的单同态。我们有一个交换图

\[
\begin{array}{c c c} \operatorname{Hom} _ {\mathbf {A}} (\mathbf {M}, \operatorname{Hom} _ {\mathbf {B}} (\mathbf {P}, \mathbf {F})) & \xrightarrow {\operatorname{Hom} _ {\mathbf {A}} (u , 1)} & \operatorname{Hom} _ {\mathbf {A}} (\mathbf {M} ^ {\prime}, \operatorname{Hom} _ {\mathbf {B}} (\mathbf {P}, \mathbf {F})) \\ \beta \Big \downarrow & & \beta^ {\prime} \Big \downarrow \\ \operatorname{Hom} _ {\mathbf {B}} (\mathbf {P} \otimes_ {\mathbf {A}} \mathbf {M}, \mathbf {F}) & \xrightarrow {\operatorname{Hom} (1 _ {\mathbf {P}} \otimes u , 1 _ {\mathbf {F}})} & \operatorname{Hom} _ {\mathbf {B}} (\mathbf {P} \otimes_ {\mathbf {A}} \mathbf {M} ^ {\prime}, \mathbf {F}) \end{array}
\]

其中 \(\beta\) 和 \(\beta'\) 是 II，第～74 页中的典范同构。由于 P 在 A 上平坦，同态 \(1_{\mathbf{P}} \otimes u: \mathbf{P} \otimes_{\mathbf{A}} \mathbf{M}' \to \mathbf{P} \otimes_{\mathbf{A}} \mathbf{M}\) 是单射。由于 F 是内射的，Hom ( \(1_{\mathbf{P}} \otimes u, 1_{\mathbf{F}}\) ) 是满射，因此 \(\operatorname{Hom}_{\mathbf{A}}(u, 1)\) 也是满射，这证明了 \(\operatorname{Hom}_{\mathbf{F}}(\mathbf{P}, \mathbf{F})\) 是内射A-模（X，第16页，引理3）。

定义3.——称A-模E为余生成元，如果对每个A-模M及M中每个非零元素x，都存在A-同态u：M → E，使得u(x) ≠ 0。

我们称A-模L为生成元，如果对于每个A-模M及M中的每个元素x，存在一个A-同态 \(u : L \to M\) 使得 \(x \in u(L)\) 。例如，A-模 \(A_{s}\) 是一个生成元。

命题12. --- 设E是一个内射A-模。要使E成为一个余生成元，必要且充分的条件是 \(\mathrm{Hom}_{\mathbf{A}}(\mathbf{S},\mathbf{E})\neq 0\) 对于每个单 A-模 S。

该条件显然是必要的。反过来，设 M 为一个 A-模，x 为 M 中的一个非零元素；M 的子模 Ax 有一个单商 S（VIII，§ 2，第 1 号，命题 3）。若 Hom \(_{A}\) (S, E) ≠ 0，则 Hom \(_{A}\) (Ax, E) ≠ 0，且存在同态 f : Ax → E 使得 f(x) ≠ 0；由于 E 是内射的，f 可延拓为从 M 到 E 的同态 u，并且我们有 u(x) = f(x) ≠ 0。

例. --- 内射 Z-模 Q/Z（X，第～18页，例2）是一个余生成元。事实上，每个单 Z-模都同构于一个模 Z/pZ， \(p \neq 0\) ，且 \(\operatorname{Hom}_{\mathbf{Z}}(\mathbf{Z}/p\mathbf{Z}, \mathbf{Q}/\mathbf{Z})\) 是非零的（例如，它包含由从 Z 到 Q 的过渡所诱导的映射 \(x \mapsto x/p\) ）。

命题13.——设B是一个环，F是一个内射余生成元B-模，P是一个(B, A)-双模。假设P在A上是平坦的，并且对每个单A-模S都有P ⊗A S ≠ 0（即，在AC，I的意义上，在A上是忠实平坦的）。则A-模Hom\_B (P, F) 是一个余生成元，并且是内射的。

事实上， \(\mathrm{Hom}_{\mathbf{B}}(\mathbf{P},\mathbf{F})\) 由命题11可知是内射的。另一方面，对于每个单A-模S， \(\mathrm{Hom}_{\mathbf{A}}(\mathbf{S},\mathrm{Hom}_{\mathbf{B}}(\mathbf{P},\mathbf{F}))\) 同构于 \(\mathrm{Hom}_{\mathbf{B}}(\mathbf{P}\otimes_{\mathbf{A}}\mathbf{S},\mathbf{F})\) ，因此是非零的，因为 \(\mathbf{P}\otimes_{\mathbf{A}}\mathbf{S}\neq 0\) 且B-模F是一个余生成元；因此由命题12可知，A-模 \(\mathrm{Hom}_{\mathbf{B}}(\mathbf{P},\mathbf{F})\) 是一个余生成元。

推论1. --- A-模 \(E_{A} = \operatorname{Hom}_{Z}(A, Q/Z)\) 是内射的且为余生成元。

我们应用命题13，取 \(\mathbf{B} = \mathbf{Z}\) , \(\mathbf{F} = \mathbf{Q} / \mathbf{Z}\) （示例）以及 \(\mathbf{P} = \mathbf{A}_d\) 。

对每个A-模M，设

\[
\mathrm{I} ^ {0} (\mathbf {M}) = \mathrm{E} _ {\mathbf {A}} ^ {\mathrm{Hom} (\mathbf {M}, \mathbf {E} _ {\mathbf {A}})}
\]

并设 \(e_{M}: M \to I^{0}(M)\) 为将每个 \(m \in M\) 对应到元素

\[
\big (\varphi (m) \big) _ {\varphi \in \operatorname{Hom} (\mathbf {M}, \mathrm{E} _ {\mathbf {A}})} \in \mathrm{I} ^ {0} (\mathbf {M}).
\]

那么：

推论2。——A-模I \(^{0}\) (M) 是内射的，且 A-同态 \(e_{\mathbf{M}} : \mathbf{M} \to \mathbf{I}^{0}(\mathbf{M})\) 是单射的。

事实上，I \(^{0}\) (M) 是内射的，因为 E \(_{A}\) 是内射的（X，第 16 页，命题 9）；另一方面，e \(_{M}\) 是单射的，因为 E \(_{A}\) 是一个余生成元。

推论3.——每个A-模都同构于某个内射A-模的子模。

推论4. --- 对于A-模E为内射的，必要且充分条件是每个单射的 A-同态 \(f: \mathbf{E} \to \mathbf{F}\) 都拥有一个 A-线性收缩。

假设E是内射的，且令 \(f: \mathrm{E} \to \mathrm{F}\) 是一个单射的A-同态。则 \(\operatorname{Hom}_{\mathbf{A}}(f, 1_{\mathbf{E}}): \operatorname{Hom}_{\mathbf{A}}(\mathrm{F}, \mathrm{E}) \to \operatorname{Hom}_{\mathbf{A}}(\mathrm{E}, \mathrm{E})\) 是满射；因此存在 \(r \in \operatorname{Hom}_{\mathbf{A}}(\mathrm{F}, \mathrm{E})\) 使得 \(r \circ f = 1_{\mathbf{E}}\) ，且 \(r\) 是 \(f\) 的 A-线性收缩。反之，存在一个内射A-模I和一个单射A-同态 \(f: \mathrm{E} \to \mathrm{I}\) (推论2)；若 \(f\) 具有一个A-线性收缩，则根据X第16页命题9，E是内射的。

